\section{Scientific Research and Student Support}\label{sec:scientific-student}

QuarkyLab is a genuinely multi-tenant platform: it carries a production neutrino-physics ML
workload, a shared teaching environment for computer-science students, and supporting AI tooling,
all on one RTX 8000. This section separates what works as a home-lab convenience from what would
be required for a formally governed student and research service.

\subsection{Research workloads}

The physicist \code{fernanda} runs a production DUNE reconstruction ML workload on the RTX 8000
(live-confirmed present and idle at review time, E-02). A DUNE Agent RAG pipeline over the
\code{dunereco} codebase assists onboarding, and a second researcher builds a separate RAG
pipeline (ARID) under a scoped exemption. Research data is protected from student starvation by a
\code{refreservation=2T} floor on the researcher dataset.

\subsection{Multi-tenant GPU sharing (live-validated 2026-07-02)}

One physical Turing card (no MIG) is shared safely between research and roughly 15 to 20 students
per semester:

\begin{itemize}
\item \rh{Scheduling.} SLURM \code{gres/shard} with 8 shards allows up to 8 concurrent GPU jobs
while leaving the card in Default compute mode so the researcher runs native.
\item \rh{Hard VRAM cap.} A per-job NVIDIA MPS instance inside each student container enforces a
per-job VRAM limit (about 6 GB per shard), chosen over SLURM \code{gres/mps} to avoid forcing
exclusive-process mode on the researcher.
\item \rh{Isolation.} \code{job\_submit.lua} wraps every student job in \code{apptainer exec --nv
--network=none} with a private home, read-only shared data, and a RAM bound; \code{cgroup.conf
ConstrainDevices=yes} denies \code{/dev/nvidia*} to any job without a GPU GRES.
\item \rh{Priority.} The research partition (\code{PriorityTier=100}) preempts student jobs
(\code{PreemptMode=REQUEUE}, not suspend, so freed VRAM is actually released), and a
requeued job auto-restarts.
\item \rh{Fairness.} Multifactor fairshare with a 3-day decay favors students who have used the
GPU least; \code{studentqos} caps one running job per student; students are batch-only.
\end{itemize}

The Phase 04 validation table shows every case passing: no-GRES jobs denied the device, a student
request for the whole card rejected, VRAM out-of-memory at the expected cap, 3 then 8 concurrent
students, the researcher native at 47.8 GB, and preemption reclaiming the card.

\subsection{Reproducibility and access lifecycle}

Reproducibility rests on a curated Apptainer image (\code{base.sif}, pinned torch and TensorFlow)
run from a local copy and refreshed daily so a storage hiccup cannot block a launch. Access is
key-only SSH over a Cloudflare Tunnel with no inbound ports; onboarding is a scripted
\code{add-cluster-key.sh} that validates the key and creates a locked-down account, and a group
sshd policy denies password auth, all forwarding, and sudo. Offboarding removes the tailnet node
and the user account.

\subsection{Readiness assessment}

{\footnotesize\sansfont
\begin{tabularx}{\linewidth}{@{}p{0.30\linewidth} X@{}}
\toprule
\rh{Strength (production-grade)} & \rh{Gap for a governed service}\\
\midrule
Real device isolation via cgroups + apptainer + MPS cap & No per-student network isolation within the node\\
Research preemption and fairshare & Research and student data both single-copy (QL-R02)\\
Per-user ZFS quotas and read-only shared data & No formal data-retention or privacy policy\\
Reproducible pinned container images & Offboarding not yet automated (planned)\\
Batch-only, no interactive un-capped access & Student containers not covered by Wazuh (QL-R22)\\
\bottomrule
\end{tabularx}}

\begin{callout}[title=Verdict]
For a home lab this is exceptional: the GPU multi-tenancy is designed and validated to a standard
most small HPC groups would accept, and the research floor plus preemption model is a real
answer to the hardest problem in shared GPU teaching. The gaps are governance rather than
mechanism: a second copy of research data (QL-R02, QL-REC03), a written retention and privacy
posture, automated account lifecycle, and extending SIEM coverage to student workloads would move
this from a well-run lab to a defensible institutional service.
\end{callout}
